\chapter{AGI-TCCU Cosmology Explorer: Simulación y Exploración Automática de Universos}
\label{chap:cosmology_explorer}

El presente capítulo describe la implementación computacional del marco TCCU-6D como un \emph{laboratorio cosmológico autónomo}. El sistema explora automáticamente grandes espacios de parámetros, simula la evolución del universo desde las fluctuaciones primordiales hasta la aparición de vida e inteligencia, y compara los resultados con observaciones astronómicas reales.

\section{Objetivo del Explorador AGI-TCCU}
El propósito es encontrar los valores de los parámetros fundamentales
\begin{equation}
\Theta = (\alpha,\beta,\tau,\Omega_m,\Omega_\Lambda,\sigma)
\end{equation}
que minimizan la diferencia entre el universo simulado y el universo observado. Los parámetros incluyen coeficientes de acoplamiento del campo $\Phi$, memoria informacional, densidad de materia y energía oscura.

\section{Función de ajuste (fitness)}
Se comparan tres observables cosmológicos críticos:
\begin{enumerate}
    \item Espectro de potencia de la materia $P(k)$.
    \item Fracción de filamentos en la red cósmica.
    \item Tamaño medio de cúmulos de galaxias.
\end{enumerate}
La función de error total es:
\begin{equation}
\chi^2 = w_1 E_{Pk} + w_2 E_{\text{web}} + w_3 E_{\text{cluster}}
\end{equation}

\section{Explorador automático de universos}
El módulo central genera parámetros aleatorios y ejecuta el simulador 6D de forma iterativa.

\begin{lstlisting}[language=Python, caption=Generación de parámetros cosmológicos]
import numpy as np

def generar_parametros():
    return {
        "alpha": np.random.uniform(0.001, 0.05),
        "beta": np.random.uniform(0.1, 1.0),
        "tau": np.random.uniform(1, 40),
        "omega_m": np.random.uniform(0.2, 0.4),
        "omega_l": np.random.uniform(0.6, 0.8),
        "sigma": np.random.uniform(0.01, 0.1)
    }
\end{lstlisting}

\subsection{Simulación y Cálculo de Métricas}
Con los parámetros elegidos, el sistema evalúa el universo resultante comparándolo con valores objetivo reales (ej. fracción de filamentos $\approx 0.18$).

\begin{lstlisting}[language=Python, caption=Función de evaluación cosmológica]
def evaluar_universo(params):
    rho = simular_universo(params)
    mean_fil, frac_fil = metricas_universo(rho)
    target_fil = 0.18
    target_len = 30
    error = (frac_fil - target_fil)**2 + (mean_fil - target_len)**2
    return error
\end{lstlisting}

\section{Evolución y Formación de Estructuras}
El simulador identifica halos gravitacionales, los asocia a galaxias y genera poblaciones estelares siguiendo la función de Salpeter ($\xi(M) \propto M^{-2.35}$).

\begin{lstlisting}[language=Python, caption=Generación de estrellas y planetas]
def generar_estrellas(galaxias):
    estrellas = []
    for g in galaxias:
        n = random.randint(10, 100)
        for _ in range(n):
            estrellas.append({"galaxia": g, "masa": random.uniform(0.1, 50)})
    return estrellas
\end{lstlisting}

\section{Influencia Galáctica y Probabilidad de Vida}
Un elemento distintivo del modelo es la conexión entre la dinámica galáctica y el origen de la vida. La probabilidad de que surja vida se ve afectada periódicamente por el cruce de los brazos espirales:
\begin{equation}
P_{\text{vida}}(t) = P_0 \left(1 + \alpha \sin\left(\frac{2\pi t}{T_{\text{arm}}}\right)\right)
\end{equation}

\section{Interpretación dentro del marco TCCU}
El Explorador AGI-TCCU materializa la cadena completa de emergencia: desde las fluctuaciones informacionales iniciales hasta las civilizaciones tecnológicas. Si se encuentra un ajuste óptimo, se obtiene un fuerte respaldo computacional para la teoría de la Creación Continua Universal.

\bigskip
\noindent\textit{“El universo ya no es solo observado: es simulado, explorado y comprendido desde su propia información.”}
